\documentclass[aspectratio=169]{beamer}
\usepackage[english]{babel}
\usepackage[utf8]{inputenc}
\usepackage[T1]{fontenc}
\usepackage{lmodern}
\usepackage{listings}
\usepackage{xcolor}
\usepackage{graphicx}
\usepackage{textcomp}
\DeclareUnicodeCharacter{2014}{---}
\DeclareUnicodeCharacter{2013}{--}
\DeclareUnicodeCharacter{2212}{-}
\DeclareUnicodeCharacter{2192}{\ensuremath{\rightarrow}}
\DeclareUnicodeCharacter{00B1}{\ensuremath{\pm}}
\DeclareUnicodeCharacter{00D7}{\ensuremath{\times}}
\DeclareUnicodeCharacter{00B7}{\ensuremath{\cdot}}
\DeclareUnicodeCharacter{20AC}{EUR}
\DeclareUnicodeCharacter{2070}{\textsuperscript{0}}
\DeclareUnicodeCharacter{00B9}{\textsuperscript{1}}
\DeclareUnicodeCharacter{00B2}{\textsuperscript{2}}
\DeclareUnicodeCharacter{00B3}{\textsuperscript{3}}
\DeclareUnicodeCharacter{2074}{\textsuperscript{4}}
\DeclareUnicodeCharacter{2075}{\textsuperscript{5}}
\DeclareUnicodeCharacter{2076}{\textsuperscript{6}}
\DeclareUnicodeCharacter{2077}{\textsuperscript{7}}
\DeclareUnicodeCharacter{2078}{\textsuperscript{8}}
\DeclareUnicodeCharacter{2079}{\textsuperscript{9}}
\DeclareUnicodeCharacter{2248}{\ensuremath{\approx}}
\DeclareUnicodeCharacter{2264}{\ensuremath{\le}}

% Colors for code
\definecolor{codebg}{RGB}{245,245,245}
\definecolor{codecomment}{RGB}{0,128,0}
\definecolor{codekeyword}{RGB}{0,0,255}
\definecolor{codestring}{RGB}{163,21,21}
\definecolor{bandcolor}{RGB}{200,50,50}

\lstdefinestyle{pythonstyle}{
    language=Python,
    backgroundcolor=\color{codebg},
    basicstyle=\ttfamily\scriptsize,
    keywordstyle=\color{codekeyword},
    commentstyle=\color{codecomment},
    stringstyle=\color{codestring},
    showstringspaces=false,
    breaklines=true,
    frame=single,
    rulecolor=\color{black},
    escapeinside={(*@}{@*)},
    moredelim=**[l][\color{bandcolor}\bfseries]{\#\ ---\ NEW},
    moredelim=**[l][\color{bandcolor}\bfseries]{\#\ ---\ CHANGED},
    literate=
      {á}{{\'a}}1 {é}{{\'e}}1 {í}{{\'i}}1 {ó}{{\'o}}1 {ú}{{\'u}}1
      {Á}{{\'A}}1 {É}{{\'E}}1 {Í}{{\'I}}1 {Ó}{{\'O}}1 {Ú}{{\'U}}1
      {ñ}{{\~n}}1 {Ñ}{{\~N}}1 {ü}{{\"u}}1 {Ü}{{\"U}}1
      {¿}{{?`}}1 {¡}{{!`}}1
      {—}{{---}}1 {–}{{--}}1 {−}{{-}}1
      {±}{{\ensuremath{\pm}}}1 {×}{{\ensuremath{\times}}}1
      {→}{{\ensuremath{\rightarrow}}}1
      {°}{{\textdegree}}1
      {·}{{\ensuremath{\cdot}}}1
      {€}{{\texteuro}}1
      {≈}{{\ensuremath{\approx}}}1
      {≤}{{\ensuremath{\le}}}1
      {⁰}{{\textsuperscript{0}}}1 {¹}{{\textsuperscript{1}}}1
      {²}{{\textsuperscript{2}}}1 {³}{{\textsuperscript{3}}}1
      {⁴}{{\textsuperscript{4}}}1 {⁵}{{\textsuperscript{5}}}1
      {⁶}{{\textsuperscript{6}}}1 {⁷}{{\textsuperscript{7}}}1
      {⁸}{{\textsuperscript{8}}}1 {⁹}{{\textsuperscript{9}}}1
}

\lstset{style=pythonstyle}

% Title slide macro: \lessontitle{Lesson Number}{Title}
\newcommand{\lessontitle}[2]{
    \title{Lesson #1: #2}
    \author{}
    \institute{Tec de Monterrey}
    \begin{frame}[plain]
        \titlepage
    \end{frame}
}

% Figure frame macro: \figureframe{Title}{Image path}
\newcommand{\figureframe}[2]{
    \begin{frame}{#1}
        \begin{center}
            \includegraphics[height=0.8\textheight,width=\textwidth,keepaspectratio]{#2}
        \end{center}
    \end{frame}
}


% Listings pulled from src/ with \lstinputlisting, so the slide is the file.
\lstset{
    basicstyle=\ttfamily\fontsize{8pt}{9.6pt}\selectfont,
    breaklines=true,
    breakatwhitespace=true,
    columns=fullflexible,
    keepspaces=true,
    xleftmargin=0.3em,
    framesep=2pt,
    aboveskip=0pt,
    belowskip=0pt
}

\newcommand{\resultframe}[3]{%
    \begin{frame}{#1}
        \begin{center}
            \includegraphics[height=0.62\textheight,width=\textwidth,keepaspectratio]{#2}
        \end{center}
        \vspace{0.25em}
        {\small #3\par}
    \end{frame}%
}

\newcommand{\claimframe}[2]{%
    \begin{frame}{#1}
        {\small #2\par}
    \end{frame}%
}

\date{}

\begin{document}

\lessontitle{12}{Adaptive GA}

\section{This lesson}

\begin{frame}{This lesson}
Lesson 07 tuned three knobs and called a fixed setting a compromise. Lesson 10 supplied the TSP. Today the parameters are allowed to move.

\vspace{0.6em}
\begin{itemize}
    \item A run can sense its own stall without knowing the optimum.
    \item The advertised adaptive branch must be audited (it may never fire).
    \item Equal-budget, paired against a \emph{tuned} rival, not against the defaults.
\end{itemize}
\end{frame}

% ================= tsp_01_fixed_parameters =================
\begin{frame}{The sensor and the update}
Let \(m_t\) be the current mean and \(\bar m\) the preceding ten-generation mean:
\[
\text{improving}_t=[m_t<(1-0.001)\bar m].
\]
Probabilities are multiplied by 0.99 while improving and 1.1 while stalled,
then clipped. Removing an individual costs no evaluation; every immigrant does.
Paired two-SE summaries below are approximate descriptions, not significance
tests.
\end{frame}

\section{One run with fixed parameters}

\begin{frame}{One run with fixed parameters}
Every lesson so far has written the parameters at the top of the file as
constants; Lesson 07 then tuned them by grid search and reported which fixed
values were best. This lesson asks a different question: should they be fixed
at all? A run has an early phase, when the population is spread out and the
search needs reach, and a late phase, when the population is nearly identical
and the search needs refinement. One constant has to serve both.
\end{frame}

\resultframe{One run with fixed parameters}{../figures/tsp_01_fixed_parameters.png}{%
    Seed 1 spends 11,901 evaluations in 99 generations and returns a legal route of \textbf{57,076}, 40.8\% longer than the nearest-neighbour baseline of 40,526}

% ================= tsp_02_the_stall =================
\section{Finding the stall}

\begin{frame}{Finding the stall}
Step 1 spent its whole budget at one setting. This step asks where that budget
went. The algorithm cannot see how far it is from the optimum -- it has no
optimum to compare against -- but it can see its own recent history, and that
is enough to tell a generation that is still making progress from one that is
not. The signal is the book's: compare the population mean of this generation
with the mean of the previous ten. If it has not improved by at least a tenth
of a percent, the generation is stalled.
\end{frame}

\begin{frame}[fragile]{(1) average()}
\lstinputlisting[basicstyle=\ttfamily\fontsize{8pt}{9.6pt}\selectfont,firstline=100,lastline=109]{../src/tsp_02_the_stall.py}
\end{frame}

\begin{frame}[fragile]{(2) is\_improving()}
\lstinputlisting[basicstyle=\ttfamily\fontsize{8pt}{9.6pt}\selectfont,firstline=112,lastline=124]{../src/tsp_02_the_stall.py}
\end{frame}

\begin{frame}[fragile]{(3) stall\_trend}
\lstinputlisting[basicstyle=\ttfamily\fontsize{8pt}{9.6pt}\selectfont,firstline=145,lastline=147]{../src/tsp_02_the_stall.py}
\end{frame}

\begin{frame}[fragile]{(4) shade\_stalls()}
\lstinputlisting[basicstyle=\ttfamily\fontsize{8pt}{9.6pt}\selectfont,firstline=180,lastline=186]{../src/tsp_02_the_stall.py}
\end{frame}

\resultframe{Finding the stall}{../figures/tsp_02_the_stall.png}{%
    The signal calls \textbf{14 of 99} generations stalled, the first at generation 74, and 14\% of the budget is spent inside them; the first half of the budget buys \textbf{85\%} of the run's total improvement, the second half 15\%}

% ================= tsp_03_adaptive_probabilities =================
\section{Probabilities that move}

\begin{frame}{Probabilities that move}
The sensor from step 2 says, every generation, whether the population is still
improving. This step wires it to the two probabilities: when the run stalls,
raise them -- more recombination, more mutation, more reach; while the run
improves, ease them down towards a floor and let selection refine what it has.
The rule is the book's, and it is deliberately blunt: multiply by 1.1 when
stalled, by 0.99 when improving, and clip.
\end{frame}

\begin{frame}[fragile]{(1) adaptation\_constants}
\lstinputlisting[basicstyle=\ttfamily\fontsize{8pt}{9.6pt}\selectfont,firstline=50,lastline=56]{../src/tsp_03_adaptive_probabilities.py}
\end{frame}

\begin{frame}[fragile]{(2) adapt\_probabilities()}
\lstinputlisting[basicstyle=\ttfamily\fontsize{8pt}{9.6pt}\selectfont,firstline=128,lastline=141]{../src/tsp_03_adaptive_probabilities.py}
\end{frame}

\begin{frame}[fragile]{(3) run(adaptive)}
\lstinputlisting[basicstyle=\ttfamily\fontsize{8pt}{9.6pt}\selectfont,firstline=148,lastline=153]{../src/tsp_03_adaptive_probabilities.py}
\end{frame}

\resultframe{Probabilities that move}{../figures/tsp_03_adaptive_probabilities.png}{%
    Same seed, same budget: adaptive \textbf{48,830} against fixed \textbf{57,076} (−14.4\%). But the stalled branch fires \textbf{0 of 99} times — crossover walks 0.90 → 0.33 and mutation 0.25 → 0.09 and never comes back up. The rule is behaving as a decay schedule, not as adaptation}

% ================= tsp_04_adaptive_population =================
\section{A population that resizes itself}

\begin{frame}{A population that resizes itself}
The second half of the book's adaptive algorithm changes the population itself.
While the run improves it drops the worst individual, so each generation gets
cheaper and the budget buys more of them. When the run stalls it injects fresh
random routes -- immigrants -- so there is new material to recombine.
\end{frame}

\begin{frame}[fragile]{(1) population\_constants}
\lstinputlisting[basicstyle=\ttfamily\fontsize{8pt}{9.6pt}\selectfont,firstline=50,lastline=54]{../src/tsp_04_adaptive_population.py}
\end{frame}

\begin{frame}[fragile]{(2) resize\_population()}
\lstinputlisting[basicstyle=\ttfamily\fontsize{7pt}{8.5pt}\selectfont,firstline=140,lastline=160]{../src/tsp_04_adaptive_population.py}
\end{frame}

\begin{frame}[fragile]{(3) run(resize)}
\lstinputlisting[basicstyle=\ttfamily\fontsize{8pt}{9.6pt}\selectfont,firstline=172,lastline=176]{../src/tsp_04_adaptive_population.py}
\end{frame}

\resultframe{A population that resizes itself}{../figures/tsp_04_adaptive_population.png}{%
    Resizing alone: 35 stalls, 70 immigrants, 103 culls, population 68–120, \textbf{138} generations for the same budget and \textbf{49,374}. With the probabilities too: 41 stalls, population down to 40, 155 generations, crossover ranging 0.39–1.00 — resizing brings back the stalls step 3 had made disappear — and \textbf{52,295}, worse on this seed than resizing alone}

% ================= tsp_05_twelve_runs =================
\section{Twelve runs of each}

\begin{frame}{Twelve runs of each}
Steps 3 and 4 each showed one seed, and one seed decides nothing: Lesson 06
measured the spread of this kind of result and Lesson 07 showed a parameter
ranking flipping from seed to seed. So every regime is now run twelve times,
run i on seed i, which makes the comparison paired -- the same twelve starting
populations face all four regimes, and the difference can be read run by run.
\end{frame}

\begin{frame}[fragile]{(1) RUNS}
\lstinputlisting[basicstyle=\ttfamily\fontsize{8pt}{9.6pt}\selectfont,firstline=56,lastline=58]{../src/tsp_05_twelve_runs.py}
\end{frame}

\begin{frame}[fragile]{(2) measure()}
\lstinputlisting[basicstyle=\ttfamily\fontsize{8pt}{9.6pt}\selectfont,firstline=238,lastline=248]{../src/tsp_05_twelve_runs.py}
\end{frame}

\begin{frame}[fragile]{(3) paired\_summary()}
\lstinputlisting[basicstyle=\ttfamily\fontsize{8pt}{9.6pt}\selectfont,firstline=251,lastline=266]{../src/tsp_05_twelve_runs.py}
\end{frame}

\begin{frame}[fragile]{(4) the\_comparison}
\lstinputlisting[basicstyle=\ttfamily\fontsize{6pt}{7.2pt}\selectfont,firstline=276,lastline=334]{../src/tsp_05_twelve_runs.py}
\end{frame}

\resultframe{Twelve runs of each}{../figures/tsp_05_twelve_runs.png}{%
    Over 12 paired runs (run \(i\) on seed \(i\)): fixed \textbf{58,195}, probabilities \textbf{50,145} (−8,050 ± 1,022, wins 12/12), resize \textbf{53,555} (−4,640 ± 926, 10/12), both \textbf{52,622} (−5,573 ± 1,127, 11/12). Every regime spends 11,901–11,946 evaluations. The approximate mean-difference ± two-SE intervals exclude zero here; this is descriptive, not a significance test}

% ================= tsp_06_a_tuned_rival =================
\section{The tuned rival}

\begin{frame}{The tuned rival}
Step 5 showed all three adaptive regimes beating the fixed algorithm, the best of
them by eight thousand units of route length. But the fixed algorithm they beat was never tuned: its
crossover and mutation probabilities were inherited from Lesson 10, where they
were chosen for a different budget and never questioned. Lesson 07 spent a whole
session on the right way to choose them, and the adaptive scheme has not yet
been asked to beat that.
\end{frame}

\begin{frame}[fragile]{(1) run(rates)}
\lstinputlisting[basicstyle=\ttfamily\fontsize{8pt}{9.6pt}\selectfont,firstline=175,lastline=189]{../src/tsp_06_a_tuned_rival.py}
\end{frame}

\begin{frame}[fragile]{(2) FIXED\_SETTINGS}
\lstinputlisting[basicstyle=\ttfamily\fontsize{8pt}{9.6pt}\selectfont,firstline=63,lastline=68]{../src/tsp_06_a_tuned_rival.py}
\end{frame}

\begin{frame}[fragile]{(3) the\_verdict}
\lstinputlisting[basicstyle=\ttfamily\fontsize{6pt}{7.2pt}\selectfont,firstline=297,lastline=366]{../src/tsp_06_a_tuned_rival.py}
\end{frame}

\resultframe{The tuned rival}{../figures/tsp_06_a_tuned_rival.png}{%
    Against four fixed settings at the same budget: 0.90/0.25 → 58,195 (adaptive wins 12/12), \textbf{0.60/0.05 → 48,870} (adaptive wins 4/12, \textbf{+1,274 ± 727}), 0.40/0.15 → 49,675 (6/12, +469 ± 480), 0.20/0.30 → 49,626 (5/12, +519 ± 835). Adaptive trails 3 of 4 means; its approximate two-SE interval against the best includes zero. ``Best'' means best among four on these same seeds}

\section{Next}

\begin{frame}{Next}
Adaptation beats the untuned setting 12/12. Against the best of four fixed settings on these same twelve seeds, the approximate two-SE interval includes zero; this exploratory study does not resolve the difference.

\vspace{0.8em}
\textbf{Lesson 13 --- Improving Performance}

Lesson 02 counted 107 evaluations instead of 110 and could not say why. Lesson 13 is the technique that makes that census the truth.
\end{frame}
\end{document}
